% === Begin Label Data ===
% ID: 132
% 难度星级: 
% 标签: 
% 备注: 
% 组卷引用次数: 0
% === End  Label Data ===

\begin{problem}{2024}{G}{北京卷}{15}{数列}
设 $\{a_n\}$ 与 $\{b_n\}$ 是两个不同的无穷数列，且都不是常数列．记集合 $M=\{k\mid a_k=b_k,k\in\mathbb{N}^*\}$，给出下列四个结论：

\circled{1} 若 $\{a_n\}$ 与 $\{b_n\}$ 均为等差数列，则 $M$ 中最多有 $1$ 个元素；

\circled{2}若 $\{a_n\}$ 与 $\{b_n\}$ 均为等比数列，则 $M$ 中最多有 $2$ 个元素；

\circled{3} 若 $\{a_n\}$ 为等差数列，$\{b_n\}$ 为等比数列，则 $M$ 中最多有 $3$ 个元素；

\circled{4} 若 $\{a_n\}$ 为递增数列，$\{b_n\}$ 为递减数列，则 $M$ 中最多有 $1$ 个元素．

其中所有正确结论的序号是 \underline{\hspace{4em}}．
\end{problem}

\begin{answer}
\circled{1}\circled{3}\circled{4}
\end{answer}

\begin{solutions}
\circled{1} 设 $a_n=a+nd$，$b_n=b+nc$．则 $k\in M$ 当且仅当 $a+kd=b+kc$，即 $a-b=k(c-d)$．这是个关于 $k$ 的方程，这个关于 $k$ 的方程只可能有 $0$ 个或 $1$ 个正整数解．特别地，对 $a_n=2n-1$，$b_n=3n-2$，此时 $M=\{1\}$．故\circled{1}正确．

\circled{2} 考虑 $a_n=2^n$，$b_n=(-2)^n$，则 $\{a_n\}$ 与 $\{b_n\}$ 都是等比数列，但是
$$
M=\{2,4,6,8,\cdots\}=\{\text{全体偶数}\},
$$
故\circled{2}错误．

\circled{3}设 $a_n=a+nd$，$b_n=q^n$（不妨设 $b_1=q$）．则 $k\in M$ 当且仅当 $a+kd=q^k$．  
设 $f(x)=|q|^x-dx-a$，$g(x)=|q|^x+dx+a$．  

（ⅰ）若 $q>0$，则 $a+kd=q^k$ 当且仅当 $f(k)=0$．求导可知 $f(x)$ 最多有两段单调区间，从而最多有 $2$ 个零点．即 $M$ 的元素个数不超过 $2$．  

（ⅱ）若 $q<0$，则当 $k$ 为奇数时，$a+kd=q^k$ 等价于 $a+kd=-|q|^k$，等价于 $g(k)=0$．  
当 $k$ 为偶数时，$a+kd=q^k$ 等价于 $f(k)=0$．  
注意到，通过讨论 $q,d$ 的取值可发现，$f(x)$ 和 $g(x)$ 必有一个在 $\mathbb{R}$ 上为单调函数，另一个在 $\mathbb{R}$ 上有两段单调区间．于是，$f(x),g(x)$ 的零点个数之和不超过 $3$．  

特别地，取 $b_n=\left(-\dfrac{1}{2}\right)^n$，则 $b_1=-\dfrac{1}{2}$，$b_3=-\dfrac{1}{8}$，$b_4=\dfrac{1}{16}$．再取 $a_n=-\dfrac{1}{2}+\dfrac{3}{16}(n-1)$，则 $a_1=-\dfrac{1}{2}$，$a_3=-\dfrac{1}{8}$，$a_4=\dfrac{1}{16}$，所以 $M=\{1,3,4\}$．故\circled{3}正确．

\circled{4} 设 $c_n=a_n-b_n$，则 $\{c_n\}$ 是递增数列，$M=\{n\mid c_n=0\}$．取一个连续函数 $f(x)$ 满足 $f(n)=c_n$，使得 $f(x)$ 是单调递增函数（构造方法：连接 $(n,c_n)$ 和 $(n+1,c_{n+1})$ 的折线）．由零点存在定理，$f(x)$ 最多有一个零点，$M$ 最多有 $1$ 个元素（反例容易举，略）．故\circled{4}正确．
\end{solutions}
